\begin{abstract}
For an elementary rule $f$, the second-order automaton $F_f$ on $(\{0,1\}^2)^{\Z}$, given by $y_{t+1}=y_{t-1}\oplus f(y_t)$, is reversible, and the automorphism group of its spacetime subshift is the group of the invertible cellular automata that commute with $F_f$. We determine this group for the rules $94$ and $95$. The automaton $95$R consists of two interleaved copies of the reversible automaton known as Rule~$54$, and we show that the group of the invertible cellular automata that commute with Rule~$54$ is $\mathcal N\times\mathcal N$, where $\mathcal N$ is the group of automorphisms of the golden mean shift that preserve the number of $1$s of every finite point. The proof uses the known description of Rule~$54$ by particles that move at light speed and are delayed at their crossings, which we state and prove in the form that is needed, for all configurations. For $95$R the group is $(\mathcal N\times\mathcal N\times\Z)/\langle(\sh^{-1},\sh,-2)\rangle$. It contains every finite group, and $\langle\sh,F_{95}\rangle$ is central of infinite index in it; yet $F_{95}$ has no blocking word and no eigenfield of any speed. Thus a large automorphism group needs neither a blocking word nor linearity. The automorphisms outside $\langle\sh,F_{95}\rangle$ have radius at least $7$ (computer-assisted). For $94$ the group is $\langle\sh,F_{94}\rangle\cong\Z^2$; the proof uses the scattering of kinks. For these rules and for rule $91$ the spacetime relation is a discrete wave equation with a source, and we prove that a block map of any radius that commutes with $F_f$ is determined by its values on the fixed points of $\sh^{\pm1}F_f$. For $91$ the group is not determined; we give partial results, and show where the method fails for rule $155$.
\end{abstract}
